The goal of the target adaptation approach is to provide a simplified and intuitive alternative to more opaque control techniques commonly seen in soft-robotics applications.
The idea is that a simple or traditional method can be used to handle the base control of the robot, and that a separate module improves the performance of that controller.
It achieves this by providing the controller with adapted targets, but without altering the controller's parameters.

In simulation, it was shown that such an approach offers near instantaneous improvement for the control of a simple mass-spring damper system.

The tests on a demonstrator robot, with a design relying on Kresling origami springs, show that the method is robust to varying controller tuning and varying system dynamics.
The target adapter module tends to make passive controllers more aggressive, and vice versa, resulting in a consistent performance.
Usage of the method on a constrained operational domain shows that both pretrained and online trained adapters manage to improve the performance, with a preference for the latter.

A limitation is presented in the form of the sampling frequency at which the signal response is stored for training the module: The frequency has to be high enough to capture enough detail.
If the sampling frequency is too low \emph{and} the recorded data contains some unstable oscillations, the adapter is unable to find meaningful relationships between the state pairs and targets it uses for updating itself.
Updates on such meaningless features lead to a negative feedback loop for the adapter's performance, which it is unable to escape.

Possible solutions to the problem are methods that allow the adapter to either recall previous settings or maintain a more stable performance.
An implementation of some kind of memory, or more envolved updating techniques like \emph{learning without forgetting}~\parencite{Li2016}, could potentially allow the adapter to escape such vicious cycles.
